As notas das avaliações do Airbnb se concentram no topo da escala e dizem pouco sobre as razões da satisfação do hóspede. Este trabalho usa o texto de 995.072 avaliações de anúncios do Rio de Janeiro (Inside Airbnb) para decompor a nota por aspecto e por queixa concreta. O sentimento das sentenças é medido por dois modelos multilíngues, e a escolha do modelo principal segue uma regra fixada antes da execução do segundo. Os aspectos vêm de tópicos obtidos com BERTopic e são checados com um segundo método (LDA) e com um segundo anotador. A medição é validada com sentenças rotuladas, ponderadas pela população, e com as subnotas oficiais, e apenas limpeza, localização e custo-benefício atendem ao critério de confiança adotado. Comodidades e limpeza concentram a maior parte da variação explicada da nota global, e as queixas de instalações, limpeza e resposta do anfitrião são as mais associadas a notas menores, com ordem semelhante nos dois modelos. Um índice sintético construído a partir do texto não acrescenta poder explicativo, além da nota oficial, sobre índices de volume de avaliações. Os resultados são associações, e a aplicação proposta é um diagnóstico por anúncio, comparado ao bairro e ao tipo de acomodação.
